\section{Related Work}

\paragraph{Projection, learned dynamics, and autonomous rollout.}
POD--Galerkin methods project governing equations onto a data-adapted basis and retain an interpretable reduced vector field \citep{sirovich1987,berkooz1993,benner2015}. Their long-horizon accuracy can deteriorate when truncated interactions are not adequately represented, motivating closure and stabilization strategies. \citep{amsallem2012}. Non-intrusive and autoregressive models instead learn reduced or latent evolution laws directly from data \citep{hesthaven2018,lee2020,champion2019,park2024,pathak2018,vlachas2018}, providing flexible nonlinear dynamics but without guaranteeing reliable recursive rollout from one-step accuracy alone. Related neural-operator approaches learn mappings between function spaces \citep{lu2021,kovachki2023}. Hybrid ROMs retain physical or projection structure while learning missing or unresolved components \citep{duraisamy2019,brunton2020,dar2023,fresca2022,geelen2023}.
RAL-MoE-ROM follows this hybrid principle, but introduces conditionally sparse corrections inside each local ROM and trains them under closed-loop rollout, explicitly targeting autonomous multi-step prediction.


\paragraph{Localized representations and conditional dynamics.}
Local reduced bases, clustering strategies, localized operators, and nonlinear manifolds have been used to reduce the representational burden placed on a single global model \citep{amsallem2008,peherstorfer2014,kaiser2014,li2021cluster,diaz2024}. These approaches motivate local reduced representations over restricted regions of the parameter or state space. Localization, however, also makes the reduced dynamics representation dependent: independently constructed local models generally differ in their affine means, bases, scalings, and operators. RAL-MoE-ROM therefore treats each local chart as a complete autonomous dynamical model, or specialist, rather than assuming that reduced coefficients from different charts are directly comparable. Its cross-regime problem is consequently not merely one of basis selection, but of conditionally assembling independently evolved local dynamical systems.

\paragraph{Mixture-of-experts and the cross-chart gap.}
Classical hierarchical mixtures learn conditional expert assignments \citep{jacobs1991,jordan1994}, while sparse MoE architectures activate only a subset of experts for each input \citep{shazeer2017,fedus2022}. Many widely used MoE formulations specialize predictors or operator blocks that operate within a shared latent representation or a common output space. Recent work extends conditional specialization to physics-informed constraints, local implicit representations, neural-operator pretraining, and PDE forecasting\citep{ICLR2024_9aeda582,NEURIPS2024_b83fae17,NEURIPS2025_2d23a999,ICLR2026_754612bd}. These approaches demonstrate the value of conditional specialization, but they do not address the same cross-chart assembly problem considered here. RAL-MoE-ROM places conditional computation at two distinct levels: sparse experts operate within a single local coordinate system, while the outer routing level acts on complete local ROMs whose coordinates are constructed independently. When two local models are combined, they are evolved independently and
coupled only after reconstruction in physical space, rather than by interpolating their reduced coefficients or operators. This distinction separates the present setting from standard within-representation MoE routing.
